\section{引言：本工作不是预测竞赛，而是结构规律发现}

固态电解质中的 Na$^+$ 迁移不是由单一几何量控制。局域 Na--X 配位决定离子离开原位点需要克服的束缚，Na 位点和空位数量决定可用载流子与落脚点，Na--Na 网络和间隙拓扑决定这些位点能否连接成长程通道，骨架柔性则可能调节迁移瓶颈随热运动开启的程度。因此，从 CIF 中计算局域结构描述符并与离子电导率关联，是建立结构设计规则的一条自然路径\cite{Adams2006,Sendek2017}。

但是，跨体系汇总数据后，“预测电导率”和“发现可迁移结构规律”不再是同一个问题。NASICON、硫化物和卤化物在阴离子尺寸、极化率、键合类型、骨架单元和典型电导率范围上均存在系统差异。一个模型只要识别出材料属于哪一类，就可能获得不错的预测分数。例如，若硫化物在当前数据中通常具有较长 Na--S 距离和较高电导率，而氧化物具有较短 Na--O 距离和较低电导率，那么全数据会呈现“Na--X 距离越长，电导率越高”。这一相关可能完全来自硫化物与氧化物之间的均值差，而不是同一体系内部改变 Na--X 距离就能提高电导率。

因此，本工作的核心问题不是

\begin{quote}
“哪个描述符最能预测 $Y$？”
\end{quote}

而是

\begin{quote}
“哪个局域结构描述符 $X$ 在去除材料体系与阴离子背景 $Z$ 后，仍与离子电导率 $Y$ 保持稳定、方向明确并可跨体系复核的关系？”
\end{quote}

这里 $X$ 是从 CIF 计算的 Na 配位、Na--Na 网络、Na 浓度、空位拓扑、骨架和对称性等结构量；$Z$ 是 NASICON、sulfide、halide 及有效阴离子对比；$Y=\log_{10}(\sigma/\mathrm{S\,cm^{-1}})$。图~\ref{fig:problem} 将这一任务拆成五个相互独立的困难。Pipeline 的每一步都不是为了堆叠模型复杂度，而是为了排除一种特定的错误结论。

\begin{figure}[htbp]
\centering
\resizebox{\textwidth}{!}{\begin{tikzpicture}[
  font=\sffamily\scriptsize,
  box/.style={draw=NatureLight,rounded corners=2pt,fill=NaturePale,minimum width=31mm,minimum height=12mm,align=center},
  risk/.style={draw=NatureRed!65,rounded corners=2pt,fill=NatureRed!6,minimum width=29mm,minimum height=16mm,align=center},
  solve/.style={draw=NatureGreen!65,rounded corners=2pt,fill=NatureGreen!7,minimum width=29mm,minimum height=16mm,align=center},
  arr/.style={-{Latex[length=2mm]},thick,draw=NatureGray}
]
\node[box] (x) at (0,0) {结构描述符 $X$\\配位、网络、空位、骨架};
\node[box] (z) at (4.6,0) {体系背景 $Z$\\NASICON / sulfide / halide};
\node[box] (y) at (9.2,0) {离子电导率 $Y$\\$\log_{10}(\sigma)$};
\draw[arr] (z.west) -- (x.east);
\draw[arr] (z.east) -- (y.west);
\node[align=center,fill=white,inner sep=1pt] at (2.3,0.55) {决定典型\\结构范围};
\node[align=center,fill=white,inner sep=1pt] at (6.9,0.55) {决定典型\\电导率范围};
\draw[arr] (x.south east) to[bend right=22] (y.south west);
\node[fill=white,inner sep=1pt] at (4.6,-0.95) {希望识别的结构关系};

\node[risk] (r1) at (0,-2.6) {体系混杂\\$X$ 可能只是 $Z$ 的代理};
\node[risk] (r2) at (2.3,-4.65) {描述符冗余\\同一物理因素被重复表达};
\node[risk] (r3) at (4.6,-2.6) {小样本不稳定\\少数材料改变筛选结果};
\node[risk] (r4) at (6.9,-4.65) {公式组合爆炸\\复杂表达产生偶然高分};
\node[risk] (r5) at (9.2,-2.6) {验证域偏移\\随机划分不等于未知体系};

\node[solve] (s1) at (0,-7.0) {秩感知控制\\Ridge 残差化};
\node[solve] (s2) at (2.3,-8.95) {物理族代表\\硬容量限制};
\node[solve] (s3) at (4.6,-7.0) {子采样 Lasso\\固定噪声校准};
\node[solve] (s4) at (6.9,-8.95) {量纲与规则约束\\二/三元公式};
\node[solve] (s5) at (9.2,-7.0) {LOSO + 多策略 CV\\V1--V4 证据};

\draw[arr] (r1) -- (s1);
\draw[arr] (r2) -- (s2);
\draw[arr] (r3) -- (s3);
\draw[arr] (r4) -- (s4);
\draw[arr] (r5) -- (s5);

\node[draw=NatureBlue!70,fill=NatureBlue!6,rounded corners=3pt,minimum width=109mm,minimum height=12mm,align=center,font=\sffamily\small\bfseries] at (4.6,-11.1)
{输出：具有透明证据链、可定位到晶体结构并值得 AIMD / 实验复核的候选假设};
\end{tikzpicture}
}
\caption{\textbf{从异质固态电解质数据库中发现结构规律面临的五类障碍。} 描述符 $X$、体系背景 $Z$ 与电导率 $Y$ 之间同时存在化学分群、描述符冗余、小样本扰动、公式搜索自由度和验证域偏移。框架的目标是逐层排除这些假发现来源，而不是用更复杂的模型掩盖它们。}
\label{fig:problem}
\end{figure}

\statusbox{\textbf{文章的结论单位。} 本框架最终输出的不是“预测器”，而是进入下一轮物理验证的候选结构假设。一个候选只有在数据可审计、去混杂关联保留、稳定性超过噪声、公式物理合法、跨体系方向可复核且不确定性可见时，才有资格进入 BVSE、NEB、AIMD 或实验验证。}

\section{为什么常规材料机器学习不能直接回答这个问题}

晶体图神经网络能够从原子连接中学习高维表示，并在拥有约 $10^4$ 或更多训练样本时展现强预测能力\cite{Xie2018,Dunn2020}。随机森林和梯度提升树也擅长拟合复杂非线性关系。然而，这些方法的主要目标是最小化测试误差。当体系类别本身携带大量目标信息时，模型可以通过学习“硫化物通常更导电”获得较高精度，却不需要学习“同一体系内哪些局域结构改变与电导率共同变化”。即使使用特征重要性，也难以区分一个变量是在描述迁移机制，还是只在识别材料类别。

此外，当前样本规模约为 84 行，远小于深度结构模型通常发挥优势的数据量。Matbench 的基准结果表明，晶体图方法的优势主要在较大数据任务中显现，而数百样本级任务中，经过合理验证的传统描述符模型仍具有竞争力\cite{Dunn2020}。近期材料机器学习研究也显示，许多所谓分布外测试实际上仍处在训练表示空间覆盖范围内，容易高估对未知化学或结构域的泛化\cite{Li2025}。因此，本研究主动限制模型容量，把计算资源用于混杂审计、稳定性和域外验证，而不是追求一个难以解释的非线性预测器。

通用符号回归和 SISSO 可以从大量初始特征中生成低维解析公式，并已用于材料描述符发现\cite{Ouyang2018}。它们与本研究追求可解释公式的方向一致，但当前数据存在两个额外问题：第一，体系标签与结构描述符强耦合，若不先去混杂，符号搜索会优先放大体系差异；第二，84 个样本不足以支持包含大量对数、幂、根式和任意比值的巨大公式空间。因此，本框架没有直接采用无界 SISSO 或遗传式符号回归，而是先压缩为稳定的物理族代表，再只执行声明式、量纲可审计的二元和受限三元组合。

\begin{table}[htbp]
\centering
\caption{Pipeline 中各算法解决的科学问题、来源与选择理由。每个模块都围绕描述符 $X$、体系背景 $Z$ 和电导率 $Y$ 的关系设计。}
\label{tab:rationale}
\small
\begin{tabularx}{\textwidth}{@{}p{0.16\textwidth}p{0.20\textwidth}p{0.18\textwidth}X@{}}
\toprule
模块 & 要排除的错误 & 主要来源领域 & 在本工作中的选择逻辑 \\
\midrule
物理描述符 & 黑箱表示无法转化为 Na$^+$ 迁移假设 & 晶体化学、图论、Voronoi 几何、材料信息学 & 把 CIF 分解为局域束缚、载流子条件、网络连通、空位和骨架响应；结果可以返回结构与材料设计。 \\
秩感知控制 & system 与 anion 重复编码同一化学分群 & 线性模型、实验设计、秩诊断 & 先判断阴离子标签是否在体系标签之外增加信息，避免把冗余 dummy 当成独立控制。 \\
Ridge 残差化 & $X$--$Y$ 相关由体系均值差伪造 & 正则化回归、流行病学协变量调整、计量经济学 partialling-out & 在共线小样本中稳定估计“该体系通常具有怎样的描述符和电导率”，再比较体系内偏离。 \\
Spearman & 结构--电导率关系可能单调但非线性 & 非参数统计 & 关注排序方向，不要求描述符变化与 $\log\sigma$ 呈严格直线；比 Pearson 更符合探索性结构规律。 \\
子采样 Lasso & 单次筛选由少数材料或相关变量竞争决定 & 高维统计、基因与生物标志物筛选 & Lasso 产生明确非零选择事件；重复半样本后，以选择频率判断一个描述符是否反复出现。 \\
固定噪声校准 & 多变量搜索中随机特征也会偶然高分 & 负对照、经验零分布、置换思想 & 让无意义变量参加同一筛选，真实描述符必须超过噪声选择频率的高分位数。 \\
物理规则组合 & 任意代数组合造成组合爆炸与伪公式 & 量纲分析、物理约束符号回归 & 只允许预先审查的族邻接、方向比值和量纲相容运算；牺牲搜索空间换取可审计性。 \\
折内 Ridge CV & 全局插补/缩放泄漏测试材料信息 & 预测模型验证 & 每个训练折独立估计中位数、尺度和回归系数；Ridge 此处只作为低容量探针，检验候选是否含可复现排序信息。 \\
LOSO 与多策略 CV & 随机划分把同体系近邻放入训练与测试 & 分组验证、domain generalization & 留出完整材料体系模拟“未知化学域”；其它策略分别检验插值能力和划分敏感性。 \\
V1--V4 & 单一相关系数无法排除公式自由度、单体系驱动和高方差 & 置换检验、正交化预测、分层分析、bootstrap & 为同一公式提供四类互补证据，并对不可用证据显式标记。 \\
\bottomrule
\end{tabularx}
\end{table}
